\chapter*{To the reader}
\addcontentsline{toc}{chapter}{To the reader}
This book starts from ordinary arithmetic. If you can add, subtract, multiply, and divide whole numbers and fractions, you have enough background to begin. No university mathematics course is assumed: no calculus, no linear algebra, no programming, and no previous experience with proofs. Chapter~1 rebuilds the algebraic notation we need, and the language of definitions, theorems, and proofs is explained before it is used. We use the real number system as a familiar tool and do not construct it from formal logic. The few properties of numbers that we accept without proof are each stated openly where they are first used: the existence of square roots (Chapter~1), the fact that integers grow beyond any real number (Chapter~2), the fact that every nonempty set of natural numbers has a smallest member (Chapter~4), and the completeness of the real numbers, which is needed to talk about limits (Chapter~11).

Many students approach a mathematics course with the question ``How much do I have to memorize before I can start?'' This book is organized around a different question: \emph{How can I enter a mathematical conversation, notice what is missing from my understanding, and learn enough to take part responsibly?} An AI assistant can be a useful partner in that conversation. It can suggest examples, explanations, arguments, and code. It cannot understand the mathematics for you. Throughout the book you will practise three skills that make such a partner useful: asking for exactly the step you are missing, noticing when an answer has quietly changed an assumption, and explaining what an argument really establishes.

Three problems run through the semester. Can a group of people be given different tasks that each of them finds acceptable? Can a process of repeated approximate updates settle down to an exact solution? Can a sequence of plus and minus signs be arranged so that it has almost no resemblance to shifted copies of itself? Each problem starts with small examples that you can work out by hand. Followed far enough, they lead to Hall's matching theorem, the contraction theorem, and a concrete entry point into current research on special matrices. A fourth strand, finite probability, shows how an average can prove that something exists without telling us directly how to build it.

Definitions appear first as answers to problems, so that you can see why they are worded as they are. Proofs appear first as plans, then as complete arguments, and finally as objects you can modify and challenge. A handful of habits---choosing a representation that reveals the goal, keeping track of assumptions, looking for an obstruction, using symmetry, and controlling error---come back in chapter after chapter, so that they become familiar through use.

\section*{How to use a chapter}
Read the opening problem and try a small example of your own before looking for a solution, whether in the text or from an assistant. Then read the explanation slowly. Each ``Pause and decide'' box asks a question and then gives a check; cover the check and commit to an answer before reading it. When you read a proof, ask two questions about every line: \emph{Why is this line true?} and \emph{Why is this line here?} You understand a proof when you can answer both.

Each chapter ends with six exercises. The first four are core exercises that practise the main ideas directly. The last two are transfer exercises that ask you to use the same ideas in a changed setting or with a changed hypothesis. Complete solutions to all six appear in Appendices~\ref{app:sola}--\ref{app:solc}. Look at a solution only after a genuine attempt, and afterwards close the book and write the argument again in your own words. The open-ended investigations near the end of the book come with worked examples and clear assessment criteria, but no answer key can promise a new theorem.

\section*{Where the course ends}
By the end of the semester you should be able to reconstruct several complete proofs from memory of their main ideas, plan what background you need for an unfamiliar statement, run and check small mathematical experiments, read short proofs written in the proof assistant Lean, and carry out a carefully limited investigation with the help of an AI assistant. You will not automatically be able to read every research paper; a specialized proof can require a great deal of further study. What you will have is a method for finding out precisely which study is needed, doing it, and connecting it back to the proof you wanted to understand.

Whether an argument is correct and whether it is new are separate questions. The theorems worked out in this book are established mathematics, presented so that you can follow and rebuild them. A student investigation becomes a research contribution only after it has been compared carefully with what is already known and its evidence has been checked.
